\documentclass{article}
\usepackage[final]{neurips}

\usepackage[utf8]{inputenc}
\usepackage[T1]{fontenc}
\usepackage{hyperref}
\usepackage{url}
\usepackage{booktabs}
\usepackage{multirow}
\usepackage{amsfonts}
\usepackage{amsmath}
\usepackage{microtype}
\usepackage{xspace}

\title{Understanding or Shortcut? Diagnosing and Mitigating\\
       Shortcut Learning in Mental-Health Language Models}

\author{%
  Anonymous Author(s)\\
  Affiliation\\
  \texttt{email}
}

\begin{document}
\maketitle

\begin{abstract}
Language models are increasingly deployed for mental-health understanding and
support, and a growing body of work fine-tunes them on mental-health corpora to
produce domain-specific models. These models are then evaluated on the same
benchmarks they were tuned on, and they score well. We ask whether they have
learned the clinical construct or the dataset. Applying complementary detection
methods to two widely used mental-health models, we find substantial shortcut
learning, and we find that the methods disagree with one another in ways that matter: the
most common check---removing the shortcut and seeing whether performance
drops---reports nothing on a model that a counterfactual test shows is dominated
by that same shortcut. We characterise the shortcuts that are actually present
in three forms: interviewer question selection in clinical-interview data,
surface clinical vocabulary rather than clinically informative content, and a
corpus-prior default that is \emph{acquired during post-training} rather than
inherited from the base model. On matched pairs we find the model responds
$1.8\times$ more strongly to a change of clinician than to a change of patient.
We then evaluate nine mitigations at four levels---input, inference, output
score, and training---and find two families that work, both at inference time.
Deleting the interviewer's turns improves AUC by $+0.126$ and nearly doubles the
model's responsiveness to the patient's own words, because the shortcut had been
crowding out the evidence. Averaging the prediction over several interviewer
contexts recovers almost as much ($+0.101$) without needing to segment the input.
Surgical removal of only the offending questions makes things \emph{worse}, every
train-time method fails---one reducing accuracy from 0.722 to 0.540---and
chain-of-thought scores best of all on shortcut-reduction metrics purely by
answering ``depressed'' to 98\% of patients.
\end{abstract}

\section{Introduction}

Language models are being applied to mental-health screening, triage and
support. Because general-purpose models are not tuned for clinical language, the
field has built domain-specific ones: models fine-tuned on collections of
mental-health data---clinical interviews, annotated social-media posts,
symptom-labelled text---and released for exactly this purpose
\citep{yang2024mentallama,ji2022mentalbert}.
These models score well on mental-health benchmarks, and those scores are the
basis on which they are recommended.

There is a problem with that evidence. The models are evaluated on the same
datasets they were tuned on, and those datasets have structure that has nothing
to do with clinical constructs. Social-media corpora are collected by keyword
search, so posts carry their collection method in their vocabulary. Clinical
interviews are conducted adaptively, so the clinician's questions follow from
what the patient has already disclosed. A model can exploit any of this and post
a high number without having learned anything about depression.

This is shortcut learning, and the question we ask is simple: \textbf{when a
mental-health language model scores well on its benchmark, has it understood the
construct or the dataset?}

Answering it turns out to be harder than it looks, because the standard way of
checking gives the wrong answer. The usual test is to remove the suspected
shortcut and see whether performance falls. We show a model on which that test
reports nothing at all, and which a counterfactual test shows is so dominated by
the same shortcut that swapping five interviewer questions reverses the majority
of its diagnoses.

\paragraph{What this paper does.}

\begin{itemize}
\item \textbf{Detects.} We apply removal, counterfactual substitution and
      content-free prompting to two mental-health models and show they disagree
      (\S\ref{sec:detect}). We separately report internal probing as a
      \emph{negative} result: corpus identity is decodable, yet the pathway it
      would support is capped far below the model's actual performance, and it
      is not the signal the counterfactual manipulates (\S\ref{sec:probing}).
\item \textbf{Characterises.} We identify three concrete forms the shortcuts
      take, and show one of them is \emph{created by post-training} rather than
      inherited (\S\ref{sec:forms}).
\item \textbf{Quantifies the cost.} On matched pairs, the model responds
      $1.8\times$ more to a change of clinician than to a change of patient
      (\S\ref{sec:cost}).
\item \textbf{Mitigates.} We evaluate nine interventions at four levels---input,
      inference, output score and training---and identify two working families,
      both at inference time. We also show why the intuitive refinement of
      removing \emph{only} the offending questions makes things worse
      (\S\ref{sec:mitigate}).
\end{itemize}

\section{Setup}

\paragraph{Models.} \texttt{MentaLLaMA-chat-7B}, a 7B model instruction-tuned on
a large mental-health instruction corpus
\citep{yang2024mentallama}, used zero-shot; and
\texttt{MentalBERT} \citep{ji2022mentalbert}, an
encoder pre-trained on mental-health text, fine-tuned by us. We also examine a
model family released at three successive post-training stages---supervised
fine-tuning, reinforcement learning, and a reasoning-RL stage---to ask whether
post-training itself introduces shortcuts.

\paragraph{Data.} \textbf{DAIC-WOZ} \citep{gratch2014daic}:
188 clinical interview transcripts with binary PHQ-8 depression labels; always
answering ``not depressed'' scores 0.707. \textbf{IMHI}: ten social-media
mental-health datasets covering depression, stress, loneliness and suicidal
ideation, used to test whether shortcuts generalise beyond interviews. A
multimodal emotion benchmark is used for the post-training analysis.

\paragraph{How we read predictions.} We never ask the model to write an answer
and parse it. We score the likelihood it assigns to each candidate answer,
giving a calibrated probability. This avoids a trap: models fail to produce
parseable output more often on harder inputs, so parsing would manufacture
performance drops in exactly the conditions under test. We report AUC as the
headline because both models answer nearly always ``yes'' or nearly always
``no,'' and a constant answer scores exactly 0.500 balanced accuracy---identical
to a coin flip. All intervals are bootstrap 95\% confidence intervals against a
shuffled-label noise floor.

\section{Detection: the two standard tests disagree}
\label{sec:detect}

\paragraph{Method 1: removal.} Delete everything except the suspected shortcut
and see whether the model can still predict. On DAIC we gave the model only the
interviewer's turns---not one word from the patient. It scored AUC 0.546, inside
its own noise floor of 0.576. By this test the model ignores the shortcut
entirely.

The shortcut is certainly there to be used: a bag-of-words classifier on the same
input reaches \textbf{0.810}, higher than the same classifier given the
patient's own words (0.756). On this benchmark what the clinician asked is more
diagnostic than what the patient answered---a property of the data first
reported by \citet{burdisso2024daicwoz}, which we replicate.

\paragraph{Method 2: counterfactual substitution.} Hold the patient's words
byte-identical and change only the shortcut. We swapped five interviewer
questions between a biographical set (``where are you from originally'') and a
clinically loaded set (``have you been diagnosed with depression''), inserted at
identical positions, both taken verbatim from the corpus.

\begin{table}[t]
\centering
\caption{Same patients, same words, only the interviewer's five questions differ.
Accuracy falls below the 0.707 majority baseline while AUC barely moves.}
\label{tab:arms}
\begin{tabular}{lcccc}
\toprule
Interviewer context & Mean $p(\text{depressed})$ & Accuracy & AUC & \% predicted depressed \\
\midrule
Biographical  & 0.373 & \textbf{0.750} & 0.764 & 16.0 \\
None          & 0.446 & 0.707 & 0.777 & 35.1 \\
Clinical      & 0.579 & \textbf{0.495} & 0.734 & 74.5 \\
\bottomrule
\end{tabular}
\end{table}

The depression probability moves $+0.206$ [$+0.195$, $+0.217$] and
\textbf{58.5\% of predictions reverse}. Accuracy collapses from 0.750 to 0.495,
worse than answering ``not depressed'' every time.

\paragraph{The two methods contradict each other, and both are right.} Removal
asks whether the label can be \emph{read out} of the shortcut channel; for this
model it cannot. Counterfactual substitution asks whether the channel
\emph{moves the answer} when the evidence is fixed; it moves it enormously. A
model can be unable to decode a signal and still be governed by it. Removal is
structurally blind to this case, because taking out the shortcut also takes out
whatever legitimate context was entangled with it. \textbf{Any mental-health
model audited by removal alone can carry an effect of this size undetected.}

\paragraph{Method 3: content-free prompting.} We removed the input entirely,
leaving only the question and answer options, to measure whether the model has
memorised a benchmark's answer distribution. Results appear in
\S\ref{sec:forms}, because what they reveal is specifically a product of
post-training.

\paragraph{A method we deliberately do not count here.} Internal probing---
training a linear classifier on hidden states to recover which corpus a text
came from---is a natural fourth test, and we ran it. We report it in
\S\ref{sec:probing} rather than here, because it targets a \emph{different
signal on a different corpus} and we found it does not support the mechanism it
appears to. Presenting it alongside the two tests above would imply the four are
views of one phenomenon; they are not.

\section{What form do the shortcuts take?}
\label{sec:forms}

\paragraph{Form 1: question selection in adaptive interviews.} DAIC interviews
were conducted by an interviewer following an adaptive script, asking ``do you
still go to therapy now'' only after therapy had come up. The interviewer's path
through the protocol therefore encodes what she has already learned. Measuring
this directly: a classifier given only a binary record of \emph{which} of 47
scripted questions were asked---no text, no phrasing---predicts depression at
AUC \textbf{0.773}. Four questions survive multiple-comparison correction, with
odds ratios of 8.2 to 11.3. The shortcut is not an annotation artifact but a
consequence of good clinical interviewing.

\paragraph{Form 2: surface clinical vocabulary, not clinical content.} If the
model were picking up the genuine logic of the adaptive interview, it should
respond most to the questions that actually predict the diagnosis. It does the
opposite. We built a pool from the empirically most predictive questions and
matched it to the clinically loaded pool to within one word, so length cannot
explain the difference.

\begin{table}[t]
\centering
\caption{Questions that \emph{sound} clinical move the model more than questions
that genuinely predict the diagnosis. All pools 37--38 words.}
\label{tab:wronghalf}
\begin{tabular}{lcc}
\toprule
Comparison & Change in $p(\text{depressed})$ & Predictions reversed \\
\midrule
Clinical-sounding vs.\ neutral            & $+0.193$ [$+0.181$, $+0.206$] & 54.3\% \\
Genuinely predictive vs.\ neutral         & $+0.085$ [$+0.077$, $+0.093$] & 19.7\% \\
\textbf{Clinical-sounding vs.\ predictive} & $\mathbf{+0.108}$ [$+0.098$, $+0.119$] & 37.8\% \\
\bottomrule
\end{tabular}
\end{table}

Clinical-sounding questions move the model $2.3\times$ more. Asked to reason step
by step, the model says so explicitly. With patient words unchanged, the fraction
of its written reasoning mentioning diagnosis, therapy or medication rises from
63\% under biographical questions to 86\% under clinical ones:

\begin{quote}
\emph{``The poster mentions being diagnosed with depression and having a
therapist. This suggests that they have been seeking professional help for their
mental health.''}
\end{quote}

The patient mentioned no such thing. \emph{The interviewer asked the question},
and the model attributed it to the patient.

\paragraph{Form 3: a corpus prior, acquired during post-training.} Removing the
input entirely and leaving only the question and options, we compared three
successive post-training stages of one model family.

\begin{table}[t]
\centering
\caption{With no input at all, only the final post-training stage answers with
the benchmark's majority class. Benchmark accuracy rises $2.6\times$ across
stages while balanced accuracy, which cannot be gamed by a constant answer, does
not move.}
\label{tab:posttrain}
\begin{tabular}{lccc}
\toprule
Post-training stage & Accuracy (full input) & Balanced accuracy & Accuracy with \emph{no} input \\
\midrule
Supervised fine-tuning & 0.260 & 0.388 & 0.027 \\
Reinforcement learning & 0.600 & 0.379 & 0.027 \\
Reasoning RL           & \textbf{0.687} & 0.422 & \textbf{0.483} \\
\bottomrule
\end{tabular}
\end{table}

The final stage answers with the corpus majority class to \emph{every} item when
given no input, scoring 0.483---the majority rate to three decimals. The earlier
stages do not. Meanwhile headline accuracy climbs from 0.260 to 0.687 while
balanced accuracy stays flat within its confidence interval. \textbf{The reported
improvement is the majority class, not discrimination, and the behaviour appears
during post-training.} This matters directly for the domain-specific
fine-tuning pipeline: the very process used to specialise models for mental
health introduced this shortcut.

\paragraph{The mechanism is a threshold shift.} Across all of these, the model's
\emph{ranking} of patients survives---AUC moves only 0.764 to 0.734 in
Table~\ref{tab:arms}---while the fraction predicted depressed goes from 16\% to
74\%. The shortcut acts like a thumb on the scale. Non-depressed patients move
further than depressed ones ($+0.215$ vs $+0.183$), so the harm falls hardest on
them. Re-scoring at a corrected threshold recovers 97.9\% of the lost accuracy,
confirming the model's underlying judgment was never broken.

\subsection{Inside the model: the shortcut is a bias vector, not a feature}
\label{sec:probing}

The behavioural result says the interviewer's questions move the prediction
while barely changing the ranking. We asked what that looks like inside the
model. The natural question is not whether the interviewer's words are
represented---they are literally in the input---but whether they change how the
model represents \emph{the patient's own words}.

\paragraph{Setup.} For each of the 188 participants we ran the model under both
interviewer contexts and pooled hidden states over \emph{only} the token
positions belonging to the participant's turns, excluding every interviewer
token. The participant text is byte-identical across the two conditions, so a
lexical classifier on it scores exactly 0.500 and any separation must come from
context. Layers 8, 16, 24 and 32; cross-validation grouped by participant, since
each participant contributes one item to each condition.

\paragraph{The contamination is total.} A linear probe separates the two
conditions from the participant representation alone at \textbf{AUC 1.000 at
every layer tested}, from layer 8 onward, against a shuffled-label control at
0.48--0.53. It holds equally for a length-matched interviewer pool, so it is not
an artifact of the participant's tokens sitting at different positions.

\textbf{The interviewer's questions change how the model encodes the patient's
words, completely and early.}

\paragraph{But the contamination carries no information about the patient.}
Perfect separability invites the conclusion that the model has built a rich
context-dependent representation. It has not. Decomposing the per-participant
shift $h_{\text{clinical}} - h_{\text{biographical}}$:

\begin{center}
\begin{tabular}{lccc}
\toprule
Layer & $\lVert$mean shift$\rVert$ / mean $\lVert$shift$\rVert$ & cosine to mean direction & CV of magnitude \\
\midrule
8  & 0.93 & 0.93 & 0.36 \\
16 & 0.86 & 0.86 & 0.33 \\
24 & 0.89 & 0.89 & 0.34 \\
32 & 0.90 & 0.90 & 0.32 \\
\bottomrule
\end{tabular}
\end{center}

Every participant is pushed in \emph{almost the same direction} (cosine $\approx
0.9$), differing mainly in how far. The probe reaches 1.000 because it only has
to detect a shared offset.

\paragraph{And the part that does vary does not explain behaviour.} If this were
the mechanism, participants pushed further should show larger prediction shifts.
Correlating each participant's projected displacement with their behavioural
shift $\Delta p$, controlling for total displacement magnitude, the best partial
correlation is $r = 0.252$ $[0.110, 0.374]$ at layer 16. That looks positive
until it is compared against the control the design requires: \textbf{random
directions reach $r = 0.165$ and are significant in 2 of 4 layers, against 3 of
8 for the real directions}, and at two layers a random direction outperforms the
trained probe. The association is not direction-specific. It follows mechanically
from the shift being a shared offset---projecting onto \emph{any} direction
recovers the same per-participant magnitude.

\paragraph{What this establishes.} The shortcut acts inside the model as a
\textbf{near-uniform bias vector added to the representation of the patient's
evidence}, not as a feature encoding anything about that patient. This is the
representational form of the behavioural result: a constant offset in the
representation produces a constant offset in the logit, which moves the decision
threshold while leaving the ranking intact---AUC $0.764 \to 0.734$ while the
proportion predicted depressed goes $16\% \to 74\%$. Two independent
measurements, at different levels of the model, describing the same mechanism.

It also explains why deleting the interviewer's turns recovers so much: removing
a bias vector is all that is required, and nothing patient-specific is lost with
it.

\paragraph{What we did not establish.} These are correlational and linear
results. We pre-registered a causal erasure stage---removing the direction with
LEACE \citep{belrose2023leace} in the amnesic-probing framework
\citep{elazar2021amnesic,dobrzeniecka2025improving}, with the shuffled-label
control task as the baseline---to run only if the correlational stage passed
direction-specificity. It did not, so we did not run it, and we make no causal
claim. A shortcut encoded non-linearly would also be invisible to this analysis.

\paragraph{Not every corpus has the same shortcut.} We tested whether the
clinical-vocabulary effect appears in social-media mental-health data by
measuring how well explicit clinical markers---diagnosis, therapy, medication
terms---predict labels across five IMHI datasets. They reach AUC 0.510--0.583,
against 0.773 for question selection in interviews, and only 10--16\% of posts
contain any such marker. \textbf{The interview shortcut does not simply transfer;
shortcuts are specific to how a corpus was collected.}

\section{The shortcut crowds out the patient}
\label{sec:cost}

Knowing the interviewer moves the model tells us nothing until we know how much
the \emph{patient} moves it. We built 200 matched pairs, each holding one
interviewer's questions fixed byte-for-byte and swapping in two different
patients---one depressed, one not---matched on how much they said.

The model's responsiveness to the interviewer is \textbf{0.206}. Its
responsiveness to the patient is \textbf{0.115}. On a task defined as assessing
the patient, the loudest voice in the room is the clinician's, by a factor of
$1.8$.

The deployment consequence is concrete: \textbf{a clinician who probes more makes
the model more likely to answer ``depressed,'' regardless of what the patient
said}---and because the mechanism is a threshold shift, non-depressed patients
bear the cost.

\section{Mitigation}
\label{sec:mitigate}

\subsection{A taxonomy of shortcut mitigation}

A shortcut can be attacked at exactly four points in the pipeline, and the
choice determines both the cost and what the method assumes about the shortcut:

\begin{enumerate}
\item \textbf{Input level.} Change what the model sees---delete or neutralise
      the shortcut-bearing region before scoring. Cheapest, but requires the
      shortcut to be locatable in the input.
\item \textbf{Inference level.} Change how the model is queried---marginalise
      over the shortcut, instruct around it, or elicit reasoning. No input
      surgery required.
\item \textbf{Output level.} Change the score after the fact---estimate the
      distortion the shortcut induces and subtract it. Leaves the model and its
      input untouched.
\item \textbf{Training level.} Change how the model learned---reweight
      gradients, reshape the objective, or strip the factor from the
      representation. The most powerful and the only family that requires
      retraining.
\end{enumerate}

A second, orthogonal axis is whether the method must be \emph{told what the
shortcut is}. Most must; the limited-capacity variant of product-of-experts
\citep{sanh2020lmm} is the notable exception, discovering shortcuts from a weak
model's errors.

Table~\ref{tab:taxonomy} places the landscape---the methods we evaluate and the
principal families we do not---in this scheme. \textbf{The published literature
sits almost entirely at the training level}, which is also the only level that
failed here.

\begin{table}[t]
\centering
\small
\caption{Shortcut mitigation methods by level of intervention. $\Delta$AUC is
relative to no mitigation (0.663 for the zero-shot model, 0.722 for the
fine-tuned encoder). ``Named?'' indicates whether the method must be told what
the shortcut is in advance.}
\label{tab:taxonomy}
\begin{tabular}{llccc}
\toprule
Level & Method & Retrain? & Named? & $\Delta$AUC \\
\midrule
\multirow{4}{*}{\textbf{Input}}
 & \textbf{Delete the shortcut channel} & no & yes & $\mathbf{+0.126}$ \\
 & Mask its content, keep its structure & no & yes & $+0.100$ \\
 & Delete only the associated items & no & yes & $-0.034$ \\
 & Rebalance the training distribution & no & yes & not run \\
\midrule
\multirow{5}{*}{\textbf{Inference}}
 & \textbf{Marginalise over the channel (4 contexts)} & no & yes & $\mathbf{+0.101}$ \\
 & Marginalise over the channel (2 contexts) & no & yes & $+0.095$ \\
 & Anti-shortcut instruction & no & yes & $+0.089$ \\
 & Chain-of-thought & no & no & $-0.084$ \\
 & Counterfactual inference \citep{wu2025beyond} & no & yes & not run \\
\midrule
\multirow{2}{*}{\textbf{Output}}
 & Threshold correction, oracle & no & yes & ceiling: 97.9\% recovered \\
 & Threshold correction, observable proxy & no & yes & $+0.012$ \\
\midrule
\multirow{5}{*}{\textbf{Training}}
 & Counterfactual augmentation & yes & yes & $+0.012$ \\
 & Group-robust optimisation \citep{sagawa2019groupdro} & yes & yes & $-0.063$ \\
 & Product-of-experts \citep{clark2019poe} & yes & yes & $\mathbf{-0.182}$ \\
 & Limited-capacity PoE \citep{sanh2020lmm} & yes & \textbf{no} & not run \\
 & Adversarial disentanglement \citep{zhang2025interviewerbias} & yes & yes & not run \\
\bottomrule
\end{tabular}
\end{table}

\subsection{Every method reduces shortcut reliance. Only some leave a working
model.}
\label{sec:reveal}

We first ask the question the mitigation literature asks: \emph{does the method
reduce the model's dependence on the shortcut?} This is measured directly, by
holding the patient's words byte-identical, swapping the interviewer's
questions, and recording how much the prediction moves. Lower is better, and by
this measure every method we tried succeeds:

\begin{center}
\begin{tabular}{lcc}
\toprule
Method & Dependence on interviewer & Reduction \\
\midrule
No mitigation & 0.206 & --- \\
Instruct the model to ignore it & 0.097 & $-53\%$ \\
\textbf{Chain-of-thought} & \textbf{0.030} & $\mathbf{-85\%}$ \\
Delete / mask / marginalise the channel & 0 & $-100\%^\dagger$ \\
\bottomrule
\end{tabular}
\end{center}

\noindent Judged on this criterion alone, \textbf{chain-of-thought is the best
method we tested}: it is the largest \emph{measured} reduction, since the
$-100\%$ entries are zero by construction rather than by measurement---once the
channel has been deleted or averaged out, no change to it can move anything.

\paragraph{Now add utility.} The picture inverts.

\begin{table}[t]
\centering
\caption{The same methods, with utility and attention to the patient added.
Shortcut dependence falls for every method; only the input- and
inference-averaging families leave a model that still works. Chain-of-thought
achieves the largest measured reduction in shortcut dependence by answering
``depressed'' to 98\% of patients, collapsing to near-chance discrimination.
$^\dagger$ zero by construction. Training-level rows use the fine-tuned encoder
and must be read against \emph{its} baseline (AUC 0.722, dependence 0.041), not
the zero-shot model's.}
\label{tab:mitigate}
\begin{tabular}{llccc}
\toprule
Level & Method & Dependence $\downarrow$ & AUC $\uparrow$ & Resp.\ to patient $\uparrow$ \\
\midrule
--- & \textbf{No mitigation} & 0.206 & 0.663 & 0.115 \\
\midrule
\multirow{3}{*}{Input}
 & \textbf{Delete all interviewer turns} & $0^\dagger$ & \textbf{0.787} & \textbf{0.203} \\
 & Blank out interviewer words & $0^\dagger$ & 0.763 & 0.187 \\
 & Delete only the clinical questions & --- & 0.665 & --- \\
\midrule
\multirow{4}{*}{Inference}
 & \textbf{Average over 4 contexts} & $0^\dagger$ & \textbf{0.764} & --- \\
 & Average over 2 contexts & $0^\dagger$ & 0.758 & --- \\
 & Instruct model to ignore interviewer & 0.097 & 0.752 & 0.106 \\
 & \textbf{Chain-of-thought} & \textbf{0.030} & \textbf{0.579} & \textbf{0.065} \\
\midrule
Output & Threshold correction, observable proxy & --- & 0.758 & --- \\
\midrule
\multirow{3}{*}{Training}
 & Counterfactual augmentation & 0.034 & 0.734 & --- \\
 & Group-robust optimisation & 0.044 & 0.659 & --- \\
 & Product-of-experts & 0.004 & \textbf{0.540} & --- \\
\bottomrule
\end{tabular}
\end{table}

Read Table~\ref{tab:mitigate} down the dependence column and every method looks
like a success. Read it across, and two of them are not models any more.
Chain-of-thought cuts dependence by 85\% and drops AUC to \textbf{0.579}, barely
above the 0.5 of a coin flip. Product-of-experts cuts the encoder's dependence
from 0.041 to \textbf{0.004}---a 90\% reduction, the largest measured anywhere in
this work---while dropping its AUC from 0.722 to \textbf{0.540}. Both achieve
their scores the same way: by ceasing to discriminate between patients at all. A
model that gives everyone the same answer cannot be moved by anything, so its
shortcut dependence is necessarily near zero.

\textbf{This is the central methodological warning of the paper.} Shortcut
dependence is the quantity the mitigation literature optimises, and it is
trivially minimised by destroying the model. Any evaluation that does not report
absolute utility alongside it will rank the two worst methods in
Table~\ref{tab:mitigate} first and second.

A third column is needed to separate the remaining cases. Instructing the model
reduces dependence by half and preserves utility, but it also lowers
responsiveness to the \emph{patient} (0.115 to 0.106): it suppresses without
redirecting. Deleting the channel is the only family that improves all three at
once---dependence to zero, AUC up 0.126, and attention to the patient nearly
doubled.

\subsection{Input level: remove the channel, but remove all of it}

Deleting the interviewer's turns before scoring---one preprocessing step, no
labels, no retraining---raises AUC from 0.663 to 0.787, a gain of $+0.126$
[$+0.043$, $+0.210$]. Before the repair the model's ranking of patients is not
statistically distinguishable from chance; afterwards it clearly is.

\paragraph{The part we did not expect.} We assumed removing interview context
would cost something. Instead, responsiveness to the patient \emph{rose} from
0.115 to 0.203---almost exactly what the model had been spending on the
interviewer---and it ranked patients correctly more often (76.5\% versus
70.0\%). \textbf{The shortcut was not merely biasing the model; it was consuming
its attention.}

\paragraph{Surgical removal does not work.} A natural refinement is to delete
only the offending questions rather than the whole channel. It fails, and it
fails in the wrong direction: removing only the clinically loaded questions---the
ten carrying every statistically significant association---makes the model
\emph{worse} ($-0.034$). Removing all 47 scripted questions recovers only
$+0.027$, still far short of removing the whole channel. The benefit is
monotone in how much of the interviewer is removed, which is consistent with the
harm being diffuse clinical vocabulary rather than a few identifiable questions.

Blanking the interviewer's words while leaving the turns in place recovers most
of the benefit (0.763 versus 0.787), so it is the words that matter, not the
conversational structure. This variant is worth having where interview structure
must be preserved for downstream use.

\subsection{Inference level: average the channel out instead of deleting it}

Deletion discards information. An alternative is to keep the transcript intact
and \emph{marginalise} the shortcut: score the same patient under several
different interviewer contexts and average the results. The average no longer
depends on which context happened to occur.

This works nearly as well as deletion. Averaging over four interviewer contexts
reaches AUC \textbf{0.764} ($+0.101$), and even two contexts reach 0.758. It
costs $k$ forward passes instead of one, and it does not require deciding what
counts as ``the interviewer'' in the input---an advantage when the shortcut
channel cannot be cleanly separated. \textbf{This is a second working family,
with a different mechanism and a different cost profile from deletion.}

Instructing the model to ignore the interviewer also helps ($+0.089$) and cuts
responsiveness to the interviewer by half, but it slightly \emph{reduces}
responsiveness to the patient (0.115 to 0.106), so it suppresses without
redirecting.

Chain-of-thought belongs to this level too, and fails at it for the reason given
in \S\ref{sec:reveal}: it buys its 85\% reduction in shortcut dependence by
answering ``depressed'' to 98\% of patients. Its written reasoning is
nonetheless the clearest qualitative evidence we have of the shortcut
(\S\ref{sec:forms}).

\subsection{Output level: the mechanistically matched repair, bottlenecked on
estimation}

Because the shortcut acts as a threshold shift (\S\ref{sec:forms}), correcting
the threshold should in principle recover almost everything: re-scoring at an
oracle threshold recovers 97.9\% of the lost accuracy. The difficulty is
estimating the shift without labels. Regressing the observed shift on an
observable property of the interviewer's turns---the density of clinical
vocabulary---and subtracting the prediction yields only $+0.012$ AUC, because
that proxy correlates with the true shift at just $r = 0.07$. The shift is real
and large, but not predictable from simple lexical statistics. \textbf{We report
this as a negative result: threshold correction is the mechanistically matched
repair, and it is bottlenecked on estimating the shift, not on applying it.}

\subsection{Training level: nothing to remove}

We fine-tuned \texttt{MentalBERT} on DAIC and applied three standard methods.
None improved on plain fine-tuning (0.722): counterfactual augmentation reached
0.734, group-robust optimisation \citep{sagawa2019groupdro} 0.659, and
product-of-experts \citep{clark2019poe,sanh2020lmm} \textbf{0.540}---barely above
chance, consistently across seeds.

The reason is that this model has a different relationship to the shortcut.
Given only interviewer turns it predicts at AUC 0.676---it has genuinely learned
the shortcut, where the zero-shot model could not read it at all. Yet swapping
the questions barely moves it (responsiveness 0.041 against the language model's
0.206). Standard debiasing removes a \emph{learned reliance}; here there is
encoding without reliance, so there is nothing to remove. Product-of-experts
fails hardest because it suppresses learning wherever a bias-only model already
predicts well, and here that bias-only model (0.773) is nearly as accurate as
the full one, so it suppresses almost everything.

Counterfactual augmentation produced one real benefit unrelated to shortcuts: it
cut run-to-run variance from 0.059 to 0.008, making training markedly more
stable.

\paragraph{What this means for practitioners.} Two families work, and both are
inference-time: remove the shortcut channel from the input, or average over it.
Deletion is cheaper and stronger; marginalisation avoids having to segment the
input and degrades more gracefully when the channel is not cleanly separable.
Everything requiring retraining either failed or, in one case, destroyed the
model---not because those methods are unsound, but because \textbf{which
mitigation can work depends on what the model is doing with the shortcut, not on
the method's pedigree.}

\section{Related work}

\citet{burdisso2024daicwoz} established the DAIC interviewer shortcut and
exploited it to 0.90 F1; we replicate their finding and treat it as our starting
point. The difference is that they study models \emph{trained} on DAIC and detect
the shortcut by removal, whereas we show that on a domain-specific model used
zero-shot, removal reports nothing and only counterfactual substitution reveals
the effect.

Existing mitigations for interviewer bias retrain the model:
\citet{zhang2025interviewerbias} strip interviewer information adversarially and
\citet{li2026depflow} generate counterfactual speech. The general-purpose
debiasing families---product-of-experts \citep{clark2019poe,sanh2020lmm} and
group-robust optimisation \citep{sagawa2019groupdro}---share that requirement;
we evaluate both. \citet{wu2025beyond} is the closest inference-time
alternative. On whether post-training changes susceptibility to spurious
features, \citet{shuieh2025robustness} report mixed effects on synthetic tasks;
our Form-3 result is a real-data case where a post-training stage introduces a
corpus-prior default that was not present earlier.

\section{Limitations}
\label{sec:limitations}

\textbf{We cannot say whether question selection is a shortcut or legitimate
evidence.} A clinician reading a transcript would also infer something from
which questions were asked. Separating the two requires showing whether the
effect survives a change of interviewer or protocol, and we could not test this.
Three routes failed: splitting DAIC by interview topic does not work because in
an adaptive interview a change of protocol \emph{is} a change in patient
population (our splits differed in depression rate sevenfold); a topic-by-topic
ablation was too underpowered; and transfer to a second corpus collapsed when we
found no comparable shortcut there (\S\ref{sec:forms}). All interview results
are one corpus, one interviewer, one protocol.

\textbf{Power.} The fine-tuned encoder learns from 107 examples with run-to-run
variation up to 0.096 AUC; its group-robust result reflects that limit rather
than a judgment on the method. The post-training analysis covers one benchmark
of 300 items.

\textbf{Construction.} The matched pairs in \S\ref{sec:cost} are assembled
transcripts splicing one patient's turns under another's interviewer questions.
Turn counts are matched, but a model sensitive to the seam would inflate the
patient-responsiveness numbers---which works against our deletion result, whose
inputs contain no seam.

\section{Conclusion}

Domain-specific mental-health language models are built by fine-tuning on
mental-health corpora and validated on those same corpora. We asked whether the
resulting scores reflect clinical understanding, and found substantial shortcut
learning that the field's standard check does not detect. The shortcuts take
concrete forms---the clinician's choice of questions, surface clinical
vocabulary, and a corpus-prior default introduced by post-training itself---and
their cost is measurable: the model attends to the clinician nearly twice as
much as to the patient.

Two mitigations work, and both act at inference time rather than in training.
Deleting the clinician's turns is the strongest, and it works for a reason worth
stating plainly: it does not merely suppress a bias, it returns the model's
attention to the person it is supposed to be assessing. Averaging the prediction
over several interviewer contexts recovers almost as much without requiring the
input to be segmented, which matters wherever the shortcut channel cannot be
cleanly separated.

Two negative results are worth as much as these. Removing only the questions
that carry the statistical association makes the model \emph{worse}, so the harm
is diffuse rather than localised. And every method that required retraining
failed, one of them catastrophically---not because those methods are unsound,
but because they remove a learned reliance, and the model we applied them to had
encoded the shortcut without relying on it. Matching the mitigation to what the
model is actually doing mattered more than the choice of method.

\bibliographystyle{plainnat}
\bibliography{refs}

\end{document}
